% GENERATED FILE. Edit canonical lessons, then run npm run generate:books.
% canonical-source-hash: fnv1a64:280a5edc70758432
% canonical-lessons: KA-C156-mundina, KA-C156-innu-munde, KA-C156-bhavishya, KA-C156-nirdharisu, KA-C156-nirikshisu

\chapter{Next, From Now On, the Future, To Decide, To Expect}
\label{ch:ka-next-from-now-on-the-future-to-decide-to-expect}
\hlchaptermodality{156}

\begin{chapteropening}
\textbf{By the end of this chapter:} \emph{I can say what will happen, and when: next week, from now on, and what I have decided.}

Complete the last lesson of chapter 156: I can say what will happen, and when: next week, from now on, and what I have decided.
\end{chapteropening}

\section[\texorpdfstring{\kn{ಮುಂದಿನ} (mundina)}{ಮುಂದಿನ (mundina)}]{\kn{ಮುಂದಿನ} (mundina) — next, coming}

\label{lesson:KA-C156-mundina}



\begin{quote}
Before the new one: say the Kannada for once, then the Kannada for to spend.
\end{quote}

\subsection*{You'll want to know: \kn{ಮುಂದಿನ}}
\textbf{\kn{ಮುಂದಿನ}} — \emph{mundina} — \textquotedblleft{}next, coming\textquotedblright{}. The future for \textquotedblleft{}I\textquotedblright{} ends in \textbf{-\kn{ಉತ್ತೇನೆ}} (\emph{-uttēne}): \textbf{\kn{ಮುಂದಿನ} \kn{ವಾರ} \kn{ನಾನು} \kn{ಬೆಂಗಳೂರಿಗೆ} \kn{ಹೋಗುತ್ತೇನೆ}} — \emph{mundina vāra nānu beṅgaḷūrige hōguttēne} — \textquotedblleft{}Next week I will go to Bengaluru.\textquotedblright{}

\begin{grammarlens}[title={What you've built}]
One more word for this chapter.
\end{grammarlens}

\subsection*{Your turn}
\begin{itemize}
\raggedright
  \item \emph{Say it:} \emph{mundina}
  \item \emph{Say it:} \emph{mundina}, once more
  \item \emph{Say it:} say \emph{mundina}
  \item \emph{Recall:} say the Kannada for once, then the Kannada for to spend, then say \emph{mundina} again
\end{itemize}

\begin{tcolorbox}[breakable,colback=teal!4,colframe=teal!35!black,title={Before you move on}]
What does \kn{ಮುಂದಿನ} mean? (\textquotedblleft{}Next, coming\textquotedblright{}.) Say it once more.
\end{tcolorbox}
\section[\texorpdfstring{\kn{ಇನ್ನು} \kn{ಮುಂದೆ} (innu munde)}{ಇನ್ನು ಮುಂದೆ (innu munde)}]{\kn{ಇನ್ನು} \kn{ಮುಂದೆ} (innu munde) — from now on}

\label{lesson:KA-C156-innu-munde}



\begin{quote}
Before the new one: say the Kannada for to spend, then the Kannada for next.
\end{quote}

\subsection*{You'll want to know: \kn{ಇನ್ನು} \kn{ಮುಂದೆ}}
\textbf{\kn{ಇನ್ನು} \kn{ಮುಂದೆ}} — \emph{innu munde} — \textquotedblleft{}from now on\textquotedblright{}. \textbf{\kn{ಇನ್ನು} \kn{ಮುಂದೆ} \kn{ನಾನು} \kn{ಬೇಗ} \kn{ಏಳುತ್ತೇನೆ}} — \emph{innu munde nānu bēga ēḷuttēne} — \textquotedblleft{}From now on I will get up early.\textquotedblright{}

\begin{grammarlens}[title={What you've built}]
One more word for this chapter.
\end{grammarlens}

\subsection*{Your turn}
\begin{itemize}
\raggedright
  \item \emph{Say it:} \emph{innu munde}
  \item \emph{Say it:} \emph{innu munde}, once more
  \item \emph{Say it:} say \emph{innu munde}
  \item \emph{Recall:} say the Kannada for to spend, then the Kannada for next, then say \emph{innu munde} again
\end{itemize}

\begin{tcolorbox}[breakable,colback=teal!4,colframe=teal!35!black,title={Before you move on}]
What does \kn{ಇನ್ನು} \kn{ಮುಂದೆ} mean? (\textquotedblleft{}From now on\textquotedblright{}.) Say it once more.
\end{tcolorbox}
\section[\texorpdfstring{\kn{ಭವಿಷ್ಯ} (bhaviṣya)}{ಭವಿಷ್ಯ (bhaviṣya)}]{\kn{ಭವಿಷ್ಯ} (bhaviṣya) — the future}

\label{lesson:KA-C156-bhavishya}



\begin{quote}
Before the new one: say the Kannada for next, then the Kannada for from now on.
\end{quote}

\subsection*{You'll want to know: \kn{ಭವಿಷ್ಯ}}
\textbf{\kn{ಭವಿಷ್ಯ}} — \emph{bhaviṣya} — \textquotedblleft{}the future\textquotedblright{}. \textbf{\kn{ಭವಿಷ್ಯದಲ್ಲಿ} \kn{ನಾನು} \kn{ಶಿಕ್ಷಕನಾಗುತ್ತೇನೆ}} — \emph{bhaviṣyadalli nānu śikṣakanāguttēne} — \textquotedblleft{}In the future I will become a teacher.\textquotedblright{}

\begin{grammarlens}[title={What you've built}]
One more word for this chapter.
\end{grammarlens}

\subsection*{Your turn}
\begin{itemize}
\raggedright
  \item \emph{Say it:} \emph{bhaviṣya}
  \item \emph{Say it:} \emph{bhaviṣya}, once more
  \item \emph{Say it:} say \emph{bhaviṣya}
  \item \emph{Recall:} say the Kannada for next, then the Kannada for from now on, then say \emph{bhaviṣya} again
\end{itemize}

\begin{tcolorbox}[breakable,colback=teal!4,colframe=teal!35!black,title={Before you move on}]
What does \kn{ಭವಿಷ್ಯ} mean? (\textquotedblleft{}The future\textquotedblright{}.) Say it once more.
\end{tcolorbox}
\section[\texorpdfstring{\kn{ನಿರ್ಧರಿಸು} (nirdharisu)}{ನಿರ್ಧರಿಸು (nirdharisu)}]{\kn{ನಿರ್ಧರಿಸು} (nirdharisu) — to decide}

\label{lesson:KA-C156-nirdharisu}



\begin{quote}
Before the new one: say the Kannada for from now on, then the Kannada for the future.
\end{quote}

\subsection*{You'll want to know: \kn{ನಿರ್ಧರಿಸು}}
\textbf{\kn{ನಿರ್ಧರಿಸು}} — \emph{nirdharisu} — \textquotedblleft{}to decide\textquotedblright{}. \textbf{\kn{ನಾನು} \kn{ಹೋಗಲು} \kn{ನಿರ್ಧರಿಸಿದೆ}} — \emph{nānu hōgalu nirdhariside} — \textquotedblleft{}I decided to go.\textquotedblright{}

\begin{grammarlens}[title={What you've built}]
One more word for this chapter.
\end{grammarlens}

\subsection*{Your turn}
\begin{itemize}
\raggedright
  \item \emph{Say it:} \emph{nirdharisu}
  \item \emph{Say it:} \emph{nirdharisu}, once more
  \item \emph{Say it:} say \emph{nirdharisu}
  \item \emph{Recall:} say the Kannada for from now on, then the Kannada for the future, then say \emph{nirdharisu} again
\end{itemize}

\begin{tcolorbox}[breakable,colback=teal!4,colframe=teal!35!black,title={Before you move on}]
What does \kn{ನಿರ್ಧರಿಸು} mean? (\textquotedblleft{}To decide\textquotedblright{}.) Say it once more.
\end{tcolorbox}
\section[\texorpdfstring{\kn{ನಿರೀಕ್ಷಿಸು} (nirīkṣisu)}{ನಿರೀಕ್ಷಿಸು (nirīkṣisu)}]{\kn{ನಿರೀಕ್ಷಿಸು} (nirīkṣisu) — to expect, to wait for}

\label{lesson:KA-C156-nirikshisu}



\begin{quote}
Before the new one: say the Kannada for the future, then the Kannada for to decide.
\end{quote}

\subsection*{You'll want to know: \kn{ನಿರೀಕ್ಷಿಸು}}
\textbf{\kn{ನಿರೀಕ್ಷಿಸು}} — \emph{nirīkṣisu} — \textquotedblleft{}to expect, to wait for\textquotedblright{}. \textbf{\kn{ನಾನು} \kn{ನಿಮ್ಮ} \kn{ಉತ್ತರವನ್ನು} \kn{ನಿರೀಕ್ಷಿಸುತ್ತೇನೆ}} — \emph{nānu nimma uttaravannu nirīkṣisuttēne} — \textquotedblleft{}I will wait for your answer.\textquotedblright{}

\begin{grammarlens}[title={What you've built}]
One more word for this chapter.
\end{grammarlens}

\subsection*{Your turn}
\begin{itemize}
\raggedright
  \item \emph{Say it:} \emph{nirīkṣisu}
  \item \emph{Say it:} \emph{nirīkṣisu}, once more
  \item \emph{Say it:} say \emph{nirīkṣisu}
  \item \emph{Recall:} say the Kannada for the future, then the Kannada for to decide, then say \emph{nirīkṣisu} again
\end{itemize}

\begin{tcolorbox}[breakable,colback=teal!4,colframe=teal!35!black,title={Before you move on}]
What does \kn{ನಿರೀಕ್ಷಿಸು} mean? (\textquotedblleft{}To expect, to wait for\textquotedblright{}.) Say it once more.
\end{tcolorbox}
